\documentclass[11pt]{article}

\usepackage{amsmath,amssymb}
\usepackage{xurl}
\usepackage[hidelinks]{hyperref}
\setlength{\emergencystretch}{2em}

% Shared commands for notes in docs/paper-gaps/.
% Notation follows the LDT paper (references/ldt-paper/).

\newcommand{\Lean}{\textsc{Lean}}
\newcommand{\leanid}[1]{\textsf{#1}}
\newcommand{\ghissue}[1]{\href{https://github.com/LionSR/MIPStarRE/issues/#1}{GitHub issue \##1}}
\newcommand{\ghpr}[1]{\href{https://github.com/LionSR/MIPStarRE/pull/#1}{GitHub PR \##1}}

% --- Paper-compatible notation ---
\newcommand{\F}{\mathbb{F}}
\newcommand{\N}{\mathbb{N}}
\newcommand{\C}{\mathbb{C}}
\newcommand{\tr}{\operatorname{tr}}
\newcommand{\ot}{\otimes}
\newcommand{\E}{\mathop{\bf E\/}}
\newcommand{\bx}{\mathbf{x}}
\newcommand{\by}{\mathbf{y}}
\newcommand{\bu}{\mathbf{u}}
\newcommand{\bv}{\mathbf{v}}

% Approximation relation (paper uses \approx_\eps and \simeq_\zeta)
\newcommand{\approxeps}{\approx_{\varepsilon}}
\newcommand{\simgeq}{\simeq_{\zeta}}

% Polynomial spaces (paper uses \polyfunc, \polysub, \polymeas)
\newcommand{\polyfunc}[3]{\operatorname{Poly}(#1,#2,#3)}
\newcommand{\polysub}[3]{\operatorname{PolySub}(#1,#2,#3)}
\newcommand{\polymeas}[3]{\operatorname{PolyMeas}(#1,#2,#3)}


\title{Issue 933: Normalization of Quantum States in the LDT Formalization}
\author{}
\date{2026-04-30}

\newcommand{\calH}{\mathcal H}
\newcommand{\ev}{\operatorname{ev}}

\begin{document}
\maketitle

\section{The normalization convention in the paper}

This note concerns the use of normalized quantum states in
\cite[Section~\emph{Preliminaries}]{Ji2020LowIndividualDegree}, and the
corresponding state object in the formal development.
\footnote{The issue record is
\href{https://github.com/LionSR/MIPStarRE/issues/933}{GitHub issue \#933}.
The paper passages consulted here are in
\path{references/ldt-paper/preliminaries.tex}: the definition of consistency is
at lines 255--268, the definition of state-dependent distance is at lines
348--360, the Cauchy--Schwarz transport proposition is at lines 509--536, the
switch-sandwich proposition is at lines 748--840, and the completion proposition
is at lines 1101--1173.}
The paper writes the shared state as a vector $|\psi\rangle$. In this
notation a state is normalized: $\langle\psi,\psi\rangle=1$. This convention
is used not only in the definitions of consistency and state-dependent distance,
but also in the proofs where a complete measurement gives a scalar equal to
$1$. For example, in the proof that the two definitions of consistency agree
for measurements, the term
\[
  \mathbb E_x \sum_b \langle\psi| I\otimes B_b^x |\psi\rangle
\]
is replaced by $1$. The same normalization convention is used when a
Cauchy--Schwarz factor of the form
$\sum_a\langle\psi|C_a^2|\psi\rangle$ is bounded by $1$.

The formal state object is deliberately more general. A formal quantum state on
$\calH$ is a positive semidefinite matrix $\rho$. It does not, by itself,
assert that $\rho$ has unit trace. The separate predicate
$\rho$ is normalized means that the normalized trace
$\tau(\rho)=\operatorname{Tr}(\rho)/\dim\calH$ is equal to $1$. The formal
expectation is
\[
  \ev_\rho(X)=\operatorname{Re} \tau(\rho X).
\]
For a pure unit vector $v$, the formal density matrix is
$\dim(\calH)\,|v\rangle\langle v|$, so that
$\tau(\dim(\calH)\,|v\rangle\langle v|)=1$ and the formal expectation agrees
with the paper's bra-ket scalar.
\footnote{The corresponding declarations are \path{QuantumState},
\path{QuantumState.IsNormalized}, \path{pureDensity},
\path{PureState.toQuantumState_isNormalized}, and \path{ev}, in
\path{MIPStarRE/LDT/Basic/QuantumState.lean}, lines 14--32, 46--108, and
159--163. The identity expectation lemma is
\path{ev_one_of_isNormalized} in
\path{MIPStarRE/LDT/Basic/OperatorExpectations.lean}, lines 88--93.}
Thus the formal development separates positivity from normalization. It also
allows mixed density matrices. This is a harmless generalization of the vector
notation when the normalization predicate is supplied.

At the strategy level the paper convention is already bundled. A symmetric or
projective LDT strategy carries both its bipartite density matrix and a proof
that it is normalized.
\footnote{See the fields \path{SymStrat.isNormalized} and
\path{ProjStrat.isNormalized} in \path{MIPStarRE/LDT/Test/StrategyCore.lean},
lines 235--246 and 371--382.}
The possible mismatch is therefore not that strategies in the main theorem lack
normalization. The mismatch is that many auxiliary lemmas are stated for an
arbitrary positive semidefinite density matrix, and a paper-scale estimate is
valid for such a lemma only if normalization is added explicitly or the estimate
is rescaled by $\ev_\rho(I)$.

\section{Why the hypothesis is mathematically necessary}

Let $\rho$ be a positive semidefinite density matrix in the formal sense, and
let $c>0$. Then $c\rho$ is again a positive semidefinite density matrix, and
\[
  \ev_{c\rho}(X)=c\,\ev_\rho(X)
\]
for every operator $X$. Consequently the formal consistency defects,
state-dependent squared distances, strong self-consistency defects, and scalar
transport averages all scale linearly with $c$.

For a normalized state and an operator $0\le X\le I$, one has
$\ev_\rho(X)\le 1$. Without normalization the bound is instead
$\ev_\rho(X)\le\ev_\rho(I)=\tau(\rho)$. This distinction is invisible in the
paper because $|\psi\rangle$ is a unit vector. It is not invisible in the
formal statement. A typical Cauchy--Schwarz step gives
\[
  |\text{transport error}|
    \le \sqrt{\delta} \sqrt{\ev_\rho(I)}.
\]
The paper bound $\sqrt\delta$ follows from this only when
$\ev_\rho(I)=1$. Thus a theorem over arbitrary positive semidefinite
$\rho$, with the same state-independent right-hand side as in the paper, would
be an overgeneralization. It would prove more than the paper's proof supplies,
and it is generally false under scalar rescaling of $\rho$.

This is a genuine mathematical issue, but it is not a defect in the chosen
state parameterization. The parameterization is sound when theorems with
paper-scale constants include the normalization hypothesis. It also gives a
useful language for intermediate weighted states and trace-scaled estimates,
where forcing normalization into the base state structure would be inconvenient.

\section{Representative formal sites where normalization is used}

The following list is representative, not exhaustive. The common pattern is
that the proof uses the identity $\ev_\rho(I)=1$, or uses it after monotonicity
to turn an operator bound $X\le I$ into a scalar bound $\ev_\rho(X)\le 1$.

First, the basic scalar bounds require normalization. The diagonal mass of a
submeasurement is at most $1$ on a normalized state, and bipartite consistency
or strong self-consistency defects are at most $1$ under the same hypothesis.
\footnote{See \path{subMeas_diagMass_le_one} in
\path{MIPStarRE/LDT/Preliminaries/SwitchSandwichPrep/Core.lean}, lines
103--114; \path{qBipartiteConsDefect_le_one_of_isNormalized},
\path{bipartiteConsError_le_one_of_isProbability}, and
\path{bipartiteConsError_uniform_le_one} in \path{MIPStarRE/LDT/Test/Defs.lean},
lines 352--414; and \path{qBipartiteSSCDefect_le_one} and
\path{bipartiteSSCError_uniform_le_one} in
\path{MIPStarRE/LDT/Pasting/BridgeLemmas/Common.lean}, lines 345--377.}
These lemmas are the formal versions of the paper's repeated use of the fact
that a bounded measurement probability is at most $1$.

Second, the Cauchy--Schwarz transport lemmas matching the paper's
$\sqrt\delta$-scale estimates carry an explicit normalization hypothesis.
This includes the transfer of consistency using state-dependent distance, the
easy approximation lemma, and both left and right clauses of the inner-product
transport proposition.
\footnote{Representative declarations are \path{triangleSub} in
\path{MIPStarRE/LDT/Preliminaries/Triangles.lean}, lines 102--124;
\path{easyApproxFromApproxDelta}, \path{closenessOfIP}, and
\path{closenessOfIPAdjoint} in \path{MIPStarRE/LDT/Preliminaries/CauchySchwarz.lean},
lines 148--218; and the underlying pointwise bounds
\path{question_easyApproxFromApproxDelta}, \path{closenessOfInnerProduct_left},
and \path{closenessOfInnerProduct_right}.}
Not every distance manipulation needs normalization. For example, a purely
operator-order contraction of a state-dependent squared distance may scale
linearly on both sides. The normalization hypothesis is needed precisely when a
bounded operator is converted into the numerical constant $1$, or when that
constant is hidden in a square-root loss.

Third, the switch-sandwich family of estimates uses normalization at exactly the
places where the paper bounds a Cauchy--Schwarz factor by $1$. The formal
statement therefore assumes that the state is normalized, and the left and right
transfer lemmas pass that assumption to the per-question gap bounds.
\footnote{See \path{switchSandwich} in
\path{MIPStarRE/LDT/Preliminaries/SwitchSandwichMain/Completeness.lean}, lines
22--45; \path{switchSandwich_leftTransfer} and \path{switchSandwich_rightTransfer}
in \path{MIPStarRE/LDT/Preliminaries/SwitchSandwichMain/LeftTransfer.lean} and
\path{MIPStarRE/LDT/Preliminaries/SwitchSandwichMain/RightTransfer.lean}; and
\path{question_switchSandwich_left_gap} and \path{question_switchSandwich_middle_gap}
in \path{MIPStarRE/LDT/Preliminaries/SwitchSandwichGapBounds/Left.lean} and
\path{MIPStarRE/LDT/Preliminaries/SwitchSandwichGapBounds/Middle.lean}.}
These lemmas feed directly into the commutativity transport estimates, such as
the scalar marginalization bounds for full slices.
\footnote{For example, see \path{fullSlice_scalar_marginalize_x} and
\path{fullSlice_scalar_marginalize_y} in
\path{MIPStarRE/LDT/Commutativity/Main/Auxiliary.lean}, lines 350--413 and
460--500, and the evaluated inner-product bridge
\path{evaluatedSlice_hEval_sandwichedRight_to_linearRight}, lines 537--640.}

Fourth, the completion and projectivization chain uses normalization when the
paper writes that the total mass of a completed measurement is $1$, or when a
truncation error is stated at the paper's scale. The formal completion theorem,
the missing-mass estimate, the spectral-truncation input, and the
orthonormalization theorem all expose the normalized-state hypothesis.
\footnote{See \path{completionMissingMassBound} in
\path{MIPStarRE/LDT/Preliminaries/SelfConsistency/Core.lean}, lines 66--149;
\path{completingToMeasurement} in
\path{MIPStarRE/LDT/Preliminaries/CompletionTransfer.lean}, lines 234--254;
\path{SpectralTruncationInput} in
\path{MIPStarRE/LDT/MakingMeasurementsProjective/Statements.lean}, lines 118--125;
\path{spectralTruncateAlmostProjective} and \path{roundAlmostProjMeas} in
\path{MIPStarRE/LDT/MakingMeasurementsProjective/Projectivization.lean}, lines
379--431; and \path{orthonormalization} in
\path{MIPStarRE/LDT/MakingMeasurementsProjective/Orthonormalization.lean}, lines
587--610.}
This is the correct formal counterpart of the paper's pure-state convention.

\section{Relation with the blueprint and Lean}

The blueprint preliminaries follow the paper's vector-state convention.  They
write $|\psi\rangle$ in the definitions of consistency and state-dependent
distance, and in the statements of the transfer lemmas, without introducing a
separate predicate for normalization.
\footnote{See \path{blueprint/src/chapter/ch03_preliminaries.tex}, lines
209--265 for \path{def:simeq} and \path{def:approx_delta}, lines 274--280 for
\path{prop:triangle-sub}, lines 316--376 for \path{prop:closeness-of-ip} and
\path{prop:easy-approx-from-approx-delta}, lines 508--536 for
\path{prop:switch-sandwich}, and lines 698--730 for
\path{prop:completing-to-measurement}.}
This matches the paper as mathematical prose: a state vector is understood to be
unit normalized.  The blueprint therefore inherits the paper convention rather
than displaying an additional normalization hypothesis.

The corresponding Lean declarations make the suppressed convention explicit.
For example, the blueprint statements for the triangle transfer, inner-product
transport, switch-sandwich lemma, completion theorem, spectral truncation input,
and orthonormalization theorem all link to declarations whose signatures include
an explicit normalized-state hypothesis.  Thus the blueprint and Lean are
compatible, but they place the convention at different levels: the blueprint uses
the standard mathematical meaning of $|\psi\rangle$ as a normalized state, while
Lean records the hypothesis as data because its base state type contains all
positive semidefinite matrices.

This note supplies the missing comparison.  Future blueprint edits should either
continue to use the paper convention for $|\psi\rangle$, or state explicitly that
a density matrix \(\rho\) is normalized whenever the prose is no longer written
in vector-state notation.  In either case, a Lean declaration linked from such a
blueprint item should carry \path{IsNormalized} whenever the proof uses the
scalar identity \(\ev_\rho(I)=1\).

\section{Verdict and migration policy}

The verdict is that \href{https://github.com/LionSR/MIPStarRE/issues/933}{issue
\#933} records a real normalization gap between the paper's convention and the
base formal state structure, but not a theorem-level defect in the current
strategy statements. The paper works with normalized states. The base formal
object intentionally records only positivity, and normalization is supplied by a
separate predicate. This is a harmless parameterization when the predicate is
required at every theorem whose conclusion has a paper-scale, state-independent
constant.

The development should therefore not immediately migrate \path{QuantumState} to
bundle normalization. Such a change would have high blast radius, and it would
remove a useful type for unnormalized or weighted positive semidefinite states.
The correct policy is local: keep \path{QuantumState} as the positive cone, keep
normalization bundled in strategy structures, and require an explicit
normalization hypothesis on any lemma that proves a bound using
$\ev_\rho(I)=1$, a scalar upper bound $\le 1$, or a paper-scale
$\sqrt\delta$ transport estimate.

When an auxiliary theorem is intentionally meant to hold for unnormalized
states, its statement should show the trace dependence explicitly. For example,
a Cauchy--Schwarz transport estimate should contain the factor
$\sqrt{\ev_\rho(I)}$, or an equivalent trace factor, rather than silently using
$1$. Conversely, if the theorem is meant to match a displayed statement in
the paper, it should assume that the state is normalized. Future audits should
look especially for occurrences of $\ev_\rho(I)=1$, scalar conclusions bounded
by $1$, and square-root transport bounds stated over an arbitrary
\path{QuantumState} without a normalization hypothesis.

\bibliographystyle{alpha}
\bibliography{references}

\end{document}
